\documentclass[11pt]{article}
\usepackage[margin=1in]{geometry}
\usepackage{hyperref}
\usepackage{booktabs}
\usepackage{longtable}
\usepackage{array}
\usepackage{xcolor}
\usepackage{enumitem}
\usepackage{titlesec}

\title{Replication Report: Heo et al.\ (2021)\\
\large Functional Genomic Insights into Probiotic \textit{Bacillus siamensis} Strain B28 from Traditional Korean Fermented \textit{Kimchi}}
\author{Ollie (OpenClaw AI subagent) \\ Set: BVBRC-61 \\ Host: CherryRd}
\date{2026-07-02}

\begin{document}
\maketitle

\section*{Paper \& Verdict}
\begin{tabular}{ll}
\textbf{Paper:} & Heo S, Kim JH, Kwak MS, Jeong DW, Sung MH. \textit{Foods} 2021, 10(8):1906. \\
\textbf{DOI:} & \href{https://doi.org/10.3390/foods10081906}{10.3390/foods10081906} \\
\textbf{PMID:} & 34441683 \quad \textbf{PMCID:} PMC8394110 \\
\textbf{Open access:} & CC BY 4.0 (Europe PMC OA full-text XML) \\
\textbf{Verdict:} & \textbf{PARTIAL REPLICATION (strong)} \\
\end{tabular}

\vspace{0.5em}
\noindent The genome-derivable core of the paper reproduces cleanly and often \emph{exactly} on the actual deposited public genome (GenBank CP066219--CP066221 / RefSeq GCF\_016313165.1): the taxonomic reclassification (ANI), the complete genome architecture (chromosome/plasmid sizes, GC, tRNA/rRNA counts --- several to the base pair), the safety genotype (no enterotoxins, no acquired AMR --- verified by two independent AMR tools), and the probiotic/functional gene inventory. What is out of reach for a sequence-only replication is the paper's wet-lab evidence (enterotoxin PCR gels, disc-diffusion antibiotic phenotypes, $\beta$-hemolysis assay, antibacterial-activity plates) --- hence PARTIAL rather than full REPLICATED.

\section{Paper Summary}
Strain B28, isolated from Korean fermented \textit{kimchi} and previously called \textit{B. polyfermenticus}, is sequenced to completion and characterized as a candidate probiotic. The authors (i) reclassify B28 as \textit{Bacillus siamensis} via 16S/MLST/ANI; (ii) report a complete genome (1 chromosome + 2 plasmids); (iii) argue the strain is \emph{safe} (no enterotoxin/hemolysin/acquired-AMR genes; susceptible to 8 antibiotics; non-$\beta$-hemolytic); and (iv) catalog probiotic-relevant gene content (bacteriocins, GABA, bile-salt hydrolase, sporulation, adhesion, stress defense, plus strain-unique $\alpha$-galactosidase and triacylglycerol lipase). This maps directly onto the BV-BRC \emph{Comprehensive Genome Analysis + AMR + Similar Genome Finder} workflow.

\section{Claims}
\begin{longtable}{clllcc}
\toprule
\# & Claim & Type & Testable? & Tested? \\
\midrule
C1 & B28 = \textit{B. siamensis} by ANI ($\geq$98\% to siamensis, $\sim$94\% to others) & Taxonomy & Yes & \checkmark \\
C2 & Complete genome: chr 3,946,178 bp + plasmids 6.1/5.4 kb; GC 45.85\%; 86 tRNA, 27 rRNA & Genome stats & Yes & \checkmark \\
C3a & No \textit{B. cereus}-type enterotoxin genes (Nhe/Hbl/CytK) & Safety genomic & Yes & \checkmark \\
C3b & \textit{hlyIII} hemolysin-III-like present; strain non-$\beta$-hemolytic & Geno.+phen. & Geno.\ yes & \checkmark / \ding{55} \\
C3c & No acquired AMR; susceptible to 8 antibiotics & Geno.+phen. & Geno.\ yes & \checkmark / \ding{55} \\
C4a & Bacteriocin/antimicrobial biosynthesis present & Genomic & Yes & \checkmark(broad) \\
C4b & Probiotic gene content (BSH/GABA/GGT/subtilisin/spor.\ etc.) & Genomic & Yes & \checkmark \\
C4c & Unique $\alpha$-galactosidase + triacylglycerol lipase & Genomic & Yes & \checkmark \\
C5 & Antibacterial activity vs foodborne pathogens (Fig.\ 2) & Wet-lab & No & \ding{55} \\
\bottomrule
\end{longtable}

\section{Methods}
All data free/public; all analysis local on CherryRd.
\begin{enumerate}[leftmargin=*,noitemsep]
    \item Paper text --- Europe PMC OA full-text XML; claims + accessions extracted.
    \item Genomes --- \texttt{datasets} CLI v18.25.1 downloaded B28 (GCF\_016313165.1, resolved from nuccore CP066219 via eutils elink) plus 6 comparators (SCSIO 05746 GCA\_002850535.1, KCTC 13613$^{\rm T}$ GCA\_000262045.1, \textit{B.\ amyloliquefaciens} FS1092/RD7-7, \textit{B.\ velezensis} JJ-D34/KMU01). B28 RefSeq annotation (\texttt{protein.faa} 3,808; \texttt{genomic.gff}) also pulled.
    \item Genome stats --- \texttt{genome\_stats.py} (contig sizes, GC, N50); tRNA/rRNA counted from RefSeq GFF.
    \item ANI (C1) --- \texttt{fastANI} + \texttt{skani} (two independent algorithms).
    \item AMR / safety (C3) --- AMRFinderPlus 4.2.7 (DB 2026-03-24.1), protein+nucleotide+GFF (\texttt{-a pgap}, \texttt{--plus}); cross-checked with RGI/CARD 3.2.7 (DIAMOND, protein mode). Enterotoxin/hemolysin from proteome scan.
    \item Functional inventory (C4) --- \texttt{func\_survey.py} regex survey of all 3,808 RefSeq products + targeted greps (GABA, subtilisin/AprX, hlyIII, bacteriocin).
    \item MLST --- \texttt{mlst} 2.33.1 (pubMLST \textit{bsubtilis} scheme).
    \item LLM-judge --- free Argo \texttt{gpt-5.2} (localhost:44497), temp 0.
\end{enumerate}

\section{Results vs.\ Paper}

\subsection{C1 --- Taxonomy / ANI (reproduced)}
\begin{tabular}{lrrrc}
\toprule
Comparison & Paper ANI & fastANI (this) & skani (this) & $>$95\%? \\
\midrule
B28 vs \textit{B.\ siamensis} KCTC 13613$^{\rm T}$ & 98.61\% & 98.42\% & 98.54\% & \checkmark \\
B28 vs \textit{B.\ siamensis} SCSIO 05746 & 97.73\% & 97.55\% & 97.67\% & \checkmark \\
B28 vs \textit{B.\ velezensis} (KMU01) & $\sim$94.06--94.28\% & 94.32\% & 94.18\% & \ding{55} \\
B28 vs \textit{B.\ amyloliquefaciens} (FS1092) & $\sim$94\% & 94.21\% & 94.19\% & \ding{55} \\
\bottomrule
\end{tabular}

\vspace{0.5em}
All values within $\sim$0.2\% of the paper. Both \textit{B.\ siamensis} comparisons exceed the 95\% species boundary; both other-species comparisons fall below. \textbf{Reclassification of B28 to \textit{B.\ siamensis} is independently reproduced.}

\subsection{C2 --- Genome architecture (reproduced, several values EXACT)}
\begin{tabular}{lllc}
\toprule
Metric & Paper & This study & Match \\
\midrule
Chromosome & 3,946,178 bp & 3,946,178 bp & \checkmark exact \\
Plasmid pB2801 & 6.1 kb & 6,117 bp & \checkmark \\
Plasmid pB2802 & 5.4 kb & 5,433 bp & \checkmark \\
Contigs total & 3 & 3 & \checkmark \\
GC\% & 45.85\% & 45.85\% & \checkmark exact \\
tRNA & 86 & 86 & \checkmark exact \\
rRNA & 27 & 27 & \checkmark exact \\
CDS & 3,573 (COG) / 1,663 (SEED) & 3,831 (PGAP) & \ding{55} pipeline-dep.\ \\
\bottomrule
\end{tabular}

\vspace{0.5em}
Comparator sanity: KCTC 13613$^{\rm T}$ returned 51 contigs (paper states its genome is \emph{incomplete}); SCSIO 05746 returned 2 contigs (complete). CDS count difference is a known annotation-pipeline artifact.

\subsection{C3 --- Safety (genotype reproduced; phenotype out of reach)}
\textbf{Enterotoxins (C3a):} proteome scan for \textit{B.\ cereus} Nhe/Hbl/CytK $\rightarrow$ \textbf{0 hits (ABSENT)}, matching PCR-verified absence. \textbf{hlyIII (C3b):} RefSeq annotates hemolysin-III-like as ``hemolysin family protein'' --- PRESENT ($\times$4). $\checkmark$ genotype.

\textbf{AMR (C3c):} AMRFinderPlus 4.2.7 reports 5 hits, \emph{all} \texttt{scope=core} (intrinsic/chromosomal): satA, fosM, two $bla$ ($\beta$-lactam), Tet(L/K/45). RGI/CARD 3.2.7: 9 Strict, 0 Perfect --- vanT/vanY (van-cluster homolog fragments, not a functional van operon), qacG/qacJ (disinfectant efflux), FosBx1, tet(45), BcI (\textit{Bacillus} cephalosporinase). Neither tool finds any \emph{acquired/mobile} resistance determinant. Fully consistent with the paper's phenotypic susceptibility to all 8 antibiotics and with the paper's own Table 2. \textbf{No acquired AMR --- upheld.}

\subsection{C4 --- Probiotic / functional gene content (reproduced)}
Survey of 3,808 RefSeq products:
\begin{itemize}[leftmargin=*,noitemsep]
    \item Bile-salt hydrolase (cholylglycine hydrolase): 1 --- \checkmark
    \item GABA: permease + $\gamma$-glutamyl-$\gamma$-aminobutyrate hydrolase + 4-aminobutyrate transaminase --- \checkmark
    \item $\gamma$-glutamyltransferase: 3 --- \checkmark
    \item Subtilisin: AprX (+ S1C/HtrA) --- \checkmark
    \item Sporulation (Spo0A etc.): 128 --- \checkmark
    \item Biofilm/EPS: 11 (EpsG, BslA, poly-$\gamma$-glutamate) --- \checkmark
    \item Adhesion (flagella + fibronectin): 41 --- \checkmark
    \item Lipoteichoic acid: 7 (LtaS, Dlt operon) --- \checkmark
    \item ROS-scavenging: 18 (catalase, SOD, GPx, AhpC, Trx) --- \checkmark
    \item Bacteriocin operons: 11 (surfactin, Blp class II, lantibiotic immunity, circular bacteriocin) --- \checkmark broad; ``bacitracin'' exact name not in RefSeq vocabulary $\triangle$
    \item \textbf{Unique:} $\alpha$-galactosidase MelA --- \checkmark
    \item \textbf{Unique:} triacylglycerol lipase (lipase/esterase $\times$6) --- \checkmark
\end{itemize}

\subsection{C5 --- Antibacterial activity (out of reach)}
Fig-2 plate assay is wet-lab; genotypic proxy (bacteriocin genes) is consistent with the claimed activity but does not constitute replication.

\section{LLM-judge (free Argo gpt-5.2)}
Independent judge: \textbf{PARTIAL}. C1 = YES; C2/C3/C4 = PARTIAL (docking C2 for annotation counts [subsequently closed by exact tRNA/rRNA], and C3/C4 for un-reproducible wet-lab and specific bacteriocin naming). Final verdict aligned, upgraded to \emph{PARTIAL (strong)} on strength of exact genome architecture + dual-tool AMR concordance.

\section{Coverage}
Of the paper's testable-from-sequence claims (C1, C2, C3a-c genotype, C4a-c), \textbf{$\sim$8/9 reproduced} (one naming caveat on bacteriocin operon). Three inherently wet-lab items (enterotoxin PCR gel, disc-diffusion phenotype, antibacterial plates) are out of reach and drive the PARTIAL verdict. No claim was contradicted.

\section{Genuine Critique}
This section is a candid, self-critical review of what this replication actually demonstrates, what it cannot demonstrate, and where the reader should temper enthusiasm.

\subsection*{Strengths (do not oversell, but do not undersell)}
\begin{itemize}[leftmargin=*,noitemsep]
    \item \textbf{Exact-match genome architecture.} Chromosome length (3,946,178 bp), GC (45.85\%), tRNA (86), and rRNA (27) counts match the paper \emph{to the base pair / to the unit}. This is a strong, mechanical form of reproducibility --- the deposited GenBank/RefSeq record \emph{is} the paper's genome.
    \item \textbf{Dual-tool AMR concordance.} AMRFinderPlus and RGI/CARD independently converge on ``no acquired AMR, only intrinsic \textit{Bacillus} homologs.'' Cross-tool concordance is stronger evidence than any single-tool run.
    \item \textbf{ANI within $\sim$0.2\%.} Reclassification is robust to algorithm choice (fastANI vs skani) and to database snapshot drift.
\end{itemize}

\subsection*{Honest limitations}
\begin{itemize}[leftmargin=*,noitemsep]
    \item \textbf{Sequence-only replication cannot touch wet-lab claims.} The enterotoxin PCR gel (C3a), disc-diffusion phenotype (C3c), $\beta$-hemolysis assay (C3b phenotype), and antibacterial-activity plates (C5) are wet-lab experiments. Absence of enterotoxin genes in the assembly is consistent with the PCR-negative result but does not \emph{prove} the strain never produced them under lab conditions the paper tested.
    \item \textbf{RefSeq annotation is not the paper's annotation.} The paper used a 2021 RAST/SEED/COG stack; we used PGAP 2024+. The CDS-count divergence (3,573 vs 3,831) is a pure annotation-pipeline artifact --- it is neither a paper error nor a replication failure, and should not be interpreted as either.
    \item \textbf{``Bacitracin operon'' naming.} The paper names a specific bacteriocin operon (``bacitracin + mesentericin''). Modern RefSeq annotates the same locus with different vocabulary (surfactin, Blp class II bacteriocin, lantibiotic immunity, circular bacteriocin). The biosynthetic capacity is unambiguously present, but the specific 2021 naming was not directly reproducible without an antiSMASH run.
    \item \textbf{Modern AMR databases surface more intrinsic hits.} AMRFinderPlus/CARD surface a few more intrinsic \textit{Bacillus} chromosomal homologs (satA, fosM, van-cluster fragments, Bc cephalosporinase) than the paper's 2021 RAST-based scan. This is \emph{not} a contradiction --- all hits are core/intrinsic and the paper's own Table 2 already notes efflux/intrinsic homologs --- but a naive reader could misinterpret ``9 CARD hits'' as evidence \emph{against} the safety claim. It is not.
    \item \textbf{Single-strain scope.} This replicates one paper about one strain (B28). It does not adjudicate the broader \textit{B.\ siamensis} vs \textit{B.\ amyloliquefaciens} vs \textit{B.\ velezensis} species-boundary questions, which are actively contested and require multi-genome core-phylogeny work (see \texttt{open\_questions.json}).
    \item \textbf{No functional validation.} The presence of a BSH gene, GABA pathway, or surfactin BGC does not prove those pathways are expressed, produced at meaningful titers, or biologically active in the kimchi matrix. Genomics establishes \emph{capacity}, not \emph{function}.
\end{itemize}

\subsection*{What a real skeptic would still want}
A truly rigorous follow-up needs: (1) antiSMASH 7.x BGC prediction to reconcile the ``bacitracin'' naming; (2) a $\geq$50-genome \textit{B.\ siamensis}/\textit{amyloliquefaciens}/\textit{velezensis} core-phylogeny with dDDH to confirm the species boundary is stable; (3) an actual proteomics or transcriptomics dataset (which does not exist for B28) to show which of the catalogued genes are expressed in kimchi; (4) a kimchi-community metagenomics comparison to establish whether B28's gene inventory is unusual for kimchi \textit{Bacillus} or typical. These are follow-on studies, not gaps in this replication.

\subsection*{Bottom line}
The paper's genome-derivable conclusions hold up under independent 2026-vintage reanalysis. The wet-lab conclusions cannot be adjudicated from sequence alone --- but nothing in the sequence contradicts them. \textbf{PARTIAL} is honest: strong on what is testable, silent on what is not.

\section{Verdict}
\textbf{PARTIAL} --- strong genome-side replication; wet-lab claims out of reach for sequence-only reanalysis; no contradictions surfaced.

\end{document}
